\section{Fluid Closure Models: Transport and Chemistry} \label{sec:fluid-closures}
This section details the transport and chemistry models necessary to
close the PDEs in the fluid model.

\subsection{Chemistry}
The chemical source terms $\dot{\omega}_{\alpha}$ are determined by
specifying a chemical mechanism and associated rate parameters.  This
section gives the form of the source terms appearing in the species
transport equations for a generic mechanism as well as example
mechanisms for argon plasma.

\subsubsection{General Form}
Assume that we have a set of $N_r$ chemical reactions, each taking the
following form:
%
\begin{equation*}
\sum_{\alpha=1}^{N_s} \nu^{\prime}_{\alpha r} X_{\alpha}
\ce{<=>}
\sum_{\alpha=1}^{N_s} \nu^{\prime\prime}_{\alpha r} X_{\alpha}
\quad
\textrm{for} \,\, r = 1, \ldots, N_r.
\end{equation*}
%
Then,
%
\begin{equation*}
\dot{\omega}_{\alpha}
=
\sum_{r=1}^{N_r} (\nu^{\prime \prime}_{\alpha r} - \nu^{\prime}_{\alpha r}) G_r
\end{equation*}
%
where $M_{\alpha}$ is the molar mass of species $\alpha$ and $G_r$ is
the rate of progress of reaction $r$.  The rate of progress is given
by
%
\begin{equation*}
G_r
=
k_{f,r} \prod_{s=1}^{N_s} n_s^{\nu^{\prime}_{sr}}
-
k_{b,r} \prod_{s=1}^{N_s} n_s^{\nu^{\prime \prime}_{sr}},
\end{equation*}
%
where $k_{f,r}$ and $k_{b,r}$ are the forward and backward rate
coefficients.  These rate coefficients are functions of temperature
that must be provided as part of the mechanism.

\subsubsection{Simple argon mechanisms}
The simplest possible argon mechanism that has been used~\cite{}
includes only irreversible ionization of ground state argon:
%
\begin{equation*}
\ce{Ar + e -> Ar^+ + 2e}.
\end{equation*}
%
Denoting this reaction 0, the energy loss per electron in $\Delta E_0
= 15.7 \si{eV}$, and the rate of progress is
%
\begin{equation*}
G_0 = k_i(T_e) n_b n_e,
\end{equation*}
%
where $n_b$ is the argon number density.  Thus,
%
\begin{equation*}
\dot{\omega}_e = (2 - 1) G_0 = k_i n_b n_e.
\end{equation*}
%
Liu et al.~\cite{} uses
%
\begin{equation*}
k_i(T_e) = A \exp(-C/T_e),
\end{equation*}
%
where $A = 1.235 \times 10^{-7} \si{cm^3.s^{-1}}$ and $C =
-18.687\si{eV}$.\todo{Ref isn't clear only units for C, but has to be eV, right?}

This may be sufficient for code development purposes, but a more
complex mechanism will be required for realistic simulations.

\subsubsection{Advanced argon mechanism}
More complex argon plasma chemistry models are detailed in~\cite{}.

\subsection{Transport}
To close the fluid model transport equation, we require transport
properties, namely the diffusion coefficient and thermal conductivity
for all species and the mobility for charged species.

\subsubsection{Simple}
Some authors~\cite{} specify these quantities as constant for a given
background number density.  Data from~\cite{} is shown in
Table~\ref{tbl:simpleTransport}.
%
\begin{table}[htp]
\caption{Transport property values from~\cite{}.}
\begin{center}
\begin{tabular}{|c|ccc|}
\hline
Name & Symbol & Units & Value \\
\hline
Electron (\ce{e}) diffusivity & $n_{\ce{Ar}} D_{\ce{e}}$ & \si{(cm.s)^{-1}} & $3.86 \times 10^{22}$ \\
Positive ion (\ce{Ar^+}) diffusivity & $n_{\ce{Ar}} D_{\ce{Ar^+}}$ & \si{(cm.s)^{-1}} & $2.07 \times 10^{18}$ \\
Metastable atom (\ce{Ar^*}) diffusivity & $n_{\ce{Ar}} D_{\ce{Ar^{\ast}}}$ & \si{(cm.s)^{-1}} & $2.42 \times 10^{18}$ \\
\hline
Electron mobility & $n_{\ce{Ar}} \mu_{\ce{e}}$ & \si{(V.cm.s)^{-1}} & $9.66 \times 10^{21}$ \\
Positive ion mobility & $n_{\ce{Ar}} \mu_{\ce{Ar^+}}$ & \si{(V.cm.s)^{-1}} & $4.65 \times 10^{19}$ \\
\hline
\end{tabular}
\end{center}
\label{tbl:simpleTransport}
\end{table}
%
The required thermal conductivities may be computed from the
diffusivity as follows:
%
\begin{equation*}
\kappa_{\alpha} = \frac{5}{2} n_{\alpha} k_B D_{\alpha}.
\end{equation*}
%
\todo[inline]{NB: this is from Panneer Chelvam and Raja.  Some authors use 3/2 instead of 5/2.  Why?  Is it an error or a different modeling choice?}


\subsubsection{Advanced}
